\begin{table}[htbp]
\centering
\caption{Construccion de la muestra de analisis}
\label{tab:sample-construction}
\begin{tabular}{lrrrr}
\toprule
Filtro & $N$ & Excluidas & \% del paso previo & \% del crudo \\
\midrule
Datos crudos (10 chunks) & 32.177 & --- & --- & 100.00 \\
Ocupados (ocu = 1) & 16.682 & 15.495 & 48.16 & 51.84 \\
Edad 18 o mas & 16.542 & 140 & 0.84 & 51.41 \\
Ingreso laboral observado & 14.764 & 1.778 & 10.75 & 45.88 \\
Horas semanales hasta 112 & 14.752 & 12 & 0.08 & 45.85 \\
Nivel educativo observado & 14.751 & 1 & 0.01 & 45.84 \\
\bottomrule
\end{tabular}
\begin{minipage}{0.95\textwidth}
\footnotesize
\textit{Notas: GEIH 2018, Bogota. Los ingresos laborales no se recortan por arriba: no hay top-coding, winsorizacion ni recorte por percentil. El ejercicio del problem set es la deteccion de subreporte de ingresos por parte de una autoridad tributaria, de modo que la cola alta es la poblacion de interes y eliminarla sesgaria justamente el objeto de estudio. Tampoco se aplica un piso de ingreso en la muestra base; el argumento income_floor permite reestimar las especificaciones principales excluyendo salarios horarios implausibles como chequeo de robustez. La ausencia de ingreso se excluye y nunca se imputa: y_total_m no registra ceros exactos, solo valores faltantes. No se aplica el factor de expansion fex_c, que se conserva como columna para los chequeos ponderados.}
\end{minipage}
\end{table}
